\documentclass[11pt]{article}

\usepackage[margin=1in]{geometry}
\usepackage{amsmath, amssymb}
\usepackage{xcolor}
\usepackage{listings}
\usepackage{hyperref}
\usepackage{array}
\usepackage{booktabs}

\hypersetup{
    colorlinks=true,
    linkcolor=blue!60!black,
    urlcolor=blue!60!black,
    citecolor=blue!60!black
}

\lstdefinestyle{bugpython}{
  language=Python,
  basicstyle=\ttfamily\small,
  keywordstyle=\color{blue!60!black},
  stringstyle=\color{green!40!black},
  commentstyle=\color{gray!70!black},
  showstringspaces=false,
  breaklines=true,
  frame=single,
  rulecolor=\color{gray!30},
  columns=fullflexible,
  keepspaces=true
}

\title{SymPy Bug Report}
\author{}
\date{}

\begin{document}

\maketitle

\section{Bug 33: solveset Represents an Unsatisfiable Tangent Equation by Infinities}

\subsection{Minimal Reproducer}

\medskip

\begin{lstlisting}[style=bugpython]
from sympy import Eq, I, S, solveset, symbols, tan

x = symbols("x")
sol = solveset(Eq(tan(x), I), x, domain=S.Complexes)
print("solution =", sol)
print("contains finite 0?", sol.contains(0))
\end{lstlisting}

\subsection{Observed Behavior}

\medskip

\begin{lstlisting}[style=bugpython]
solution = ImageSet(Lambda(_n, _n*pi + oo*I), Integers)
contains finite 0? False
\end{lstlisting}

SymPy returns a periodic image set whose elements are of the form
\(\_n\pi+\infty i\). This is not a set of finite complex solutions to the
equation \(\tan(x)=i\). The recorded reproducer also evaluates the residual at
\(x=0\) as \(-1.0 i\), confirming that an ordinary finite value is not being
silently represented by the returned expression.

\subsection{Expected Behavior}

The equation \(\tan(x)=i\) has no finite complex solution, so
\[
\operatorname{solveset}(\tan(x)=i,\ x,\ \mathbb{C}) = \varnothing.
\]
The same conclusion holds for the paired singular target \(\tan(x)=-i\).

\subsection{Mathematical Explanation}

For finite complex \(z\),
\[
\tan(z) = -i\,\frac{\exp(2iz)-1}{\exp(2iz)+1}.
\]
If \(\tan(z)=i\), then
\[
-i\,\frac{\exp(2iz)-1}{\exp(2iz)+1}=i
\quad\Longrightarrow\quad
\frac{\exp(2iz)-1}{\exp(2iz)+1}=-1
\quad\Longrightarrow\quad
\exp(2iz)=0.
\]
The complex exponential is never zero at a finite complex argument, so the
equation has no finite complex solution. For \(\tan(z)=-i\), the same formula
would require an infinite exponential value instead of a finite one. Thus the
returned \(oo*I\)-valued image set is not a harmless branch choice, a missing
integer parameter, or an equivalent representation of a finite solution family;
it represents a singular inverse value as if it were a valid complex solution.

\subsection{Source-Level Diagnosis}

The diagnosis identifies the relevant solver logic in
\nolinkurl{sympy/solvers/solveset.py}, specifically the complex trigonometric
inversion helper around lines 467--473:

\begin{lstlisting}[style=bugpython]
_invert_trig_hyp_complex
\end{lstlisting}

The traced call path starts with \texttt{solveset}, which rewrites the equation
as \(\tan(x)-i=0\) in \texttt{\_solveset}. The trigonometric equation detector
then routes the expression to \texttt{\_solve\_trig}. Because the equation
contains one trigonometric atom, \texttt{\_solve\_trig} calls \texttt{\_invert};
the complex inversion path moves the additive \(-i\) and inverts
\(\tan(x)=i\).

At that point, the same helper applies tangent's periodic inverse pattern,
effectively constructing \(n\pi+\arctan(a)\) for a finite
target \(a\). For \(a=i\), however, \(\arctan(i)\) evaluates to \(oo*I\) through
the pure-imaginary branch in
\nolinkurl{sympy/functions/elementary/trigonometric.py}. The solver then embeds
that non-finite inverse value into an \texttt{ImageSet} and returns it without a
later finite-complex filter or an empty-preimage check.

This source-level behavior causes the mathematical failure because tangent's
singular target values \(i\) and \(-i\) do not have finite complex preimages,
but the inversion rule treats the infinite principal inverse as though it were a
valid seed for the periodic solution family. The supported fix direction is to
detect non-finite inverse values in the tangent and cotangent complex inverse
branches, at least for \(a=\pm i\), and return \texttt{EmptySet}; more generally,
inverse image sets should be filtered to finite complex values.

\subsection{Additional Failing Instances}

The independently checked failure family consists of the two singular tangent
targets \(i\) and \(-i\). These are the values where the inverse tangent is
non-finite rather than a finite complex seed for the usual \(n\pi+\arctan(a)\)
solution family. Both cases should therefore return the empty set over
\(\mathbb{C}\).

\begin{center}
\small
\renewcommand{\arraystretch}{1.3}
\begin{tabular}{@{}>{\raggedright\arraybackslash}p{0.44\linewidth}>{\raggedright\arraybackslash}p{0.23\linewidth}>{\raggedright\arraybackslash}p{0.23\linewidth}@{}}
\toprule
\textbf{Input / Expression} & \textbf{SymPy behavior} & \textbf{Expected behavior} \\
\midrule
\(\operatorname{solveset}(\tan(x)=i,\ x,\ \mathbb{C})\) &
\(\operatorname{ImageSet}(\Lambda(\_n,\ \_n\pi+oo\,i),\ \mathbb{Z})\) &
\(\varnothing\) \\
\(\operatorname{solveset}(\tan(x)=-i,\ x,\ \mathbb{C})\) &
not \texttt{EmptySet}; contains a non-finite inverse tangent value &
\(\varnothing\) \\
\bottomrule
\end{tabular}
\end{center}

The expected behavior is the same across the family because neither equation
can be satisfied by finite complex \(x\): one would force
\(\exp(2ix)=0\), and the other would require the corresponding exponential term
to become infinite.

\subsection{Related Issues Assessment}

A best-effort deduplication check found public material related only at the
general level of trigonometric inversion and \texttt{ImageSet} representations
for infinite families. It did not identify a public issue, pull request, Stack
Overflow post, or release-note entry for \(\tan(x)=i\) returning an
\(oo*I\)-valued image set. The deduplication verdict is therefore that this is a
new instance of a known failure mode: complex trigonometric inversion must
reject non-finite inverse images, but this exact reproducible SymPy case appears
previously unreported and remains individually actionable as an issue and a
regression test.

\subsection{Proposed SymPy Regression Test}

\medskip

\begin{lstlisting}[style=bugpython]
import pytest
from sympy import Eq, I, S, solveset, symbols, tan

x = symbols("x")
CASES = [I, -I]

@pytest.mark.parametrize("target", CASES)
def test_tan_singular_targets_have_no_finite_complex_solution(target):
    assert solveset(Eq(tan(x), target), x, domain=S.Complexes) is S.EmptySet
\end{lstlisting}

\end{document}
